% Generated draft: requires mathematical and source review.
\documentclass[11pt]{article}
\usepackage[margin=25mm]{geometry}
\usepackage{fontspec}
\setmainfont{Latin Modern Roman}
\usepackage{amsmath,amssymb,mathtools,unicode-math}
\setmathfont{Latin Modern Math}
\usepackage{enumitem,needspace,xurl,hyperref}
\setlength{\parindent}{0pt}
\setlength{\parskip}{5pt}
\setlength{\emergencystretch}{3em}
\newcommand{\N}{\mathbb{N}}
\newcommand{\Z}{\mathbb{Z}}
\newcommand{\Q}{\mathbb{Q}}
\newcommand{\R}{\mathbb{R}}
\newcommand{\C}{\mathbb{C}}
\newcommand{\F}{\mathbb{F}}
\newcommand{\A}{\mathbb{A}}
\newcommand{\PP}{\mathbb{P}}
\newcommand{\E}{\mathbb{E}}
\newcommand{\eps}{\varepsilon}
\newcommand{\dd}{\,\mathrm{d}}
\begin{document}
\section*{UnsolvedMath ID: 800765}
\textbf{Source problem code:} catalogue entry
\subsection{Definitions and mathematical statement}
Let \(K\) be an arbitrary field and let \(t\) be an indeterminate over \(K\). A finite split embedding problem over \(K\) consists of a finite Galois extension \(L/K\), a finite group \(G\), and a surjective homomorphism
\[
\alpha:G\twoheadrightarrow\operatorname{Gal}(L/K)
\]
that admits a homomorphic section \(s:\operatorname{Gal}(L/K)\to G\), so that \(\alpha\circ s=\mathrm{id}\). Identify \(\operatorname{Gal}(L(t)/K(t))\) with \(\operatorname{Gal}(L/K)\) by restriction to \(L\).

Using the standard proper-solution interpretation, a regular solution consists of a finite Galois extension \(F/K(t)\) containing \(L(t)\) and an isomorphism
\[
\gamma:\operatorname{Gal}(F/K(t))\xrightarrow{\sim}G
\]
such that
\[
\alpha\circ\gamma
=\operatorname{res}_{L(t)/K(t)},
\]
where the restriction map on the right is viewed as taking values in \(\operatorname{Gal}(L/K)\), and such that \(F/L\) is regular. Here regular means that \(F/L\) is separably generated and every element of \(F\) algebraic over \(L\) belongs to \(L\).

\textbf{Conjecture (D\`ebes--Deschamps, as attributed in the supplied source).} For every field \(K\), every finite Galois extension \(L/K\), every finite group \(G\), and every split surjective homomorphism \(\alpha:G\twoheadrightarrow\operatorname{Gal}(L/K)\), there exists a regular solution as above.

The supplied excerpt calls \(\gamma\) a solution without explicitly specifying its surjectivity or bijectivity. The formulation above uses the conventional proper-solution meaning; this interpretation requires source review. The excerpt also says that this conjecture combines two preceding conjectures, whose statements are not included in the supplied record.
\subsection{Short English statement}
The conjecture asserts that every finite split embedding problem over every field can be realized by a Galois extension of a rational function field, with the prescribed finite Galois extension as its field of constants. The statement uses the proper-solution interpretation described above. Its open status is supplied registry information.
\subsection{Sources}
\begin{itemize}
\item Local UnsolvedMath registry, permanent ID \(800765\), source code \texttt{catalogue entry}; supplied upstream revision \texttt{d371ee917ac9647fbb06cd6fc529deec6f3dda82}. The registry labels the problem \emph{open}; this is a supplied classification, not an independently verified literature assessment.
\item Bary-Soroker and Fehm, \emph{Open Problems in the Theory of Ample Fields} (2011), Conjecture 8. Supplied source URL: \url{https://arxiv.org/abs/1106.1310}. The record reports extraction from arXiv LaTeX and access on 2026-10-16; this access information is supplied provenance, not a claim of independent consultation.
\item The supplied statement and background attribute the conjecture to D\`ebes and Deschamps. No separate original reference for that attribution is supplied.
\end{itemize}
\end{document}
